\section{The Jurisprudential Core: The Migration of the Anchor of Validity}
\label{p13:sec:jurisprudence}

\noindent
The diagnosis ended by sending the inquiry to the philosophy of law, and for a definite reason. Law is the institution modernity charged with enforceability; it is where the device that the diagnosis found failing\,---\,force as the producer of expectation\,---\,was housed, theorised, and refined. If that device is not, after all, the cause of obligation, then the law's own reflection upon itself should bear the marks of the discovery, and the history of legal philosophy should record, in its internal development, a movement away from the very ground modernity had assigned it. This section shows that it does. It is not a survey. It reads the major positions in the philosophy of law along a single axis\,---\,\emc{the location of the anchor of validity}, the place at which each position lodges the force by which an obligation obliges\,---\,and it finds, reading them so, not a set of competing doctrines to be weighed but a directional movement to be traced: the anchor migrates, across a century of the law's reflection upon itself, from force toward trust. The section follows that migration through eight positions, and arranges them at its close in a single view.

A word on the relation to the foundational paper, which treated several of these same thinkers along a different axis, that of structure against generation, the Symbolic against the Real~\parencite{wanhong_grb}. The present reading is not a repetition but a re-cutting of the same material along a more focused axis: not what kind of thing law is, but where the force of its obligation is anchored. The two readings are consistent and complementary; the present one is the one the problem of trust requires.

\subsection{The anchor in force: Austin and Hobbes}
\label{p13:sub:jur-austin}

The movement begins where modernity's assignment is most undisguised, in the command theory of John Austin. For Austin, law is the command of a sovereign, backed by the threat of a sanction, habitually obeyed; and a legal obligation is, by definition, the standing liability to the sanction\,---\,to be obliged is to be exposed to the probability of an evil at the hands of one with the power and the will to inflict it~\parencite{austin1832}. The anchor of validity here is welded to mechanical cause. The obligation \emph{is} the threatened pain; there is, in the account, no remainder, nothing in being bound over and above being liable to be hurt. Whatever else the command theory does, it states with unmatched clarity the position the rest of the chapter will move away from: that the force of an obligation is the force of a threatened consequence, and nothing else.

Behind Austin stands his metaphysical ancestor, Hobbes, whose famous sentence is the chapter's true point of departure: covenants without the sword are but words, and of no strength to secure a man at all~\parencite{hobbes1651}. The sentence is usually read as a worldly observation about the weakness of unbacked promises. It is better read, for the purposes of this paper, as the statement of an anthropology from which the impossibility of trust follows by necessity. Hobbes's natural man, in the war of all against all, has no ground on which to expect his fellow's forbearance, and so trust is not merely scarce in the state of nature but \emph{logically impossible}: there is nothing from which the expectation could be formed, and the sovereign must therefore be erected to inject, from outside, by the threat of overwhelming force, the expectation the parties cannot generate between themselves.

It is exactly here that the chapter's reading first bites. Austin and Hobbes take the limiting case\,---\,the case in which trust has wholly collapsed and only fear remains\,---\,and mistake it for the whole. The Hobbesian state of nature is the crisis of the symbol carried to its extremity, a world in which no symbol can produce expectation because none can be relied upon; and the Hobbesian remedy, the sword, is precisely the substitution of mechanical cause for the inferential cause that has failed. The command theory, then, is not false so much as it is the theory of a pathological limit. Its model has nothing to say where trust is present, because there force is redundant; it becomes visible, and seems complete, only where trust has collapsed and force is all that is left. And this is why it speaks so directly to the present: the contemporary overload of justice (\xref{p13:sub:diag-justice}) is the world sliding toward the Hobbesian condition\,---\,trust thinning until the sword is asked to do everything\,---\,at the very moment when, as the symptomatology showed, the sword can no longer fall. The command theory describes the destination toward which the failure of trust tends; and the failure of the sword to fall, there, is the command theory's own solution failing on its own terms.

\subsection{The anchor moves toward acceptance: Hart's internal point of view}
\label{p13:sub:jur-hart}

The decisive turn, and the hinge of the whole chapter, is Hart's. His critique of Austin is the moment at which the anchor of validity is, for the first time within legal positivism itself, pried loose from mechanical cause and moved toward something else. Two of his objections do the work. The first is the distinction between being \emph{obliged} and having an \emph{obligation}~\parencite{hart1961}. The gunman obliges me to hand over my money; I am not under an obligation to do so. Austin's theory, which defines obligation as liability to a sanction, can capture only the first\,---\,the being-compelled\,---\,and is constitutively unable to capture the second, the normativity by which a rule gives me a reason and not merely a motive. The second objection is that whole classes of law\,---\,the power-conferring rules by which one makes a will, forms a contract, marries\,---\,are not commands backed by threats at all, but enabling structures that confer the capacity to alter one's normative situation, and that the command theory simply cannot accommodate.

From these objections Hart draws the concept on which everything turns: the \emph{internal point of view}~\parencite{hart1961}. A rule exists, as a rule and not as a mere observed regularity of compelled behaviour, where there is a group who accept it from the inside, who take it as a standard for their own conduct and as a ground for criticising deviation in others\,---\,who use the rule, in their reasoning and their appraisals, rather than merely conforming to it under threat. The internal point of view is a normative attitude, not a prediction; and it stands in sharp contrast, as \xref{p13:sub:jur-holmes} will show, to the external, predictive stance of one who asks only what will be done to him if he is caught. This is the jurisprudential moment at which the anchor of validity is moved from mechanical cause toward the normative\,---\,toward an acceptance that gives reasons, the very thing the inferential cause of the diagnosis traffics in.

The point reaches its depth in Hart's account of the \emph{rule of recognition}, the ultimate rule by which a legal system's other rules are identified as valid. What grounds the rule of recognition itself? Hart's answer is that nothing grounds it from above; it rests upon nothing but the convergent practice of the officials who accept and use it. Its validity is not conferred by a higher rule, and is not secured by any sanction; its existence simply \emph{is} the fact that it is accepted and acted upon. Here Hart arrives, without naming it so, at the position closest to this paper's: the ultimate ground of a legal order's validity is not a sword but a collective, self-sustaining acceptance\,---\,a convergence of the internal point of view across the body of officials, which holds itself up by being shared. This is, in the vocabulary of \xref{p13:sec:mechanism}, the institutional-scale form of the structural channel: a shared prior, sustained by the very fact of being shared, that needs no external pump to hold it in place.

Honesty requires that Hart's own difficulty be marked, for it is the difficulty that hands the argument to Fuller. Is the acceptance in which the rule of recognition consists a matter of fact or a matter of norm? Hart wishes it to be a convergent social fact; but a bare factual convergence\,---\,that officials do, as it happens, behave alike\,---\,seems unable to generate the genuine obligation, the \emph{ought}, that a rule of recognition is supposed to carry. This is the opening through which Dworkin pressed his critique, that law cannot be a matter of convergent practice alone but must involve an interpretive, value-laden commitment~\parencite{dworkin1986}. The paper does not adjudicate that dispute here; it notes only that the difficulty points beyond bare acceptance toward something acceptance alone does not supply, and that Fuller named what that something is.

\subsection{The anchor in reciprocity: Fuller's inner morality of law}
\label{p13:sub:jur-fuller}

Fuller supplies what bare acceptance lacks. His claim is that law possesses an inner morality, a set of eight principles\,---\,that rules be general, promulgated, prospective, clear, non-contradictory, possible to obey, stable over time, and administered congruently with their announced terms\,---\,which are not external moral constraints upon law but conditions of its very existence as law~\parencite{fuller1969}. The question Fuller's critics pressed is why these should be called a \emph{morality} rather than a mere technique of effective governance, a craftsman's know-how for making rules that work. His answer, and the point at which his account becomes this paper's, is that the eight principles together constitute an implicit promise from the lawgiver to the governed, and that the eighth\,---\,\emph{congruence}, the requirement that officials actually administer the law as announced\,---\,is the heart of the matter, because it makes obedience a \emph{reciprocity}. The citizen's fidelity to law is a response to the state's fidelity to its own announced rules; each side's keeping of faith is conditioned upon the other's. Strip the reciprocity away\,---\,let the state announce one rule and apply another\,---\,and the obligation to obey does not merely weaken but loses its ground, because the relation that grounded it was a mutual keeping of faith and one party has defected from it.

This reciprocity is trust under its proper name in the law. When Fuller says that law is a collaborative enterprise between the government and the citizen, an enterprise that fails not when a rule is broken but when the bond of reciprocity that constitutes it is betrayed, he has named, in the jurisprudential register, the very thing this paper is about: the anchor of validity lodged in a mutually relied-upon undertaking. The long debate between Fuller and Hart over whether law and morality are necessarily connected can, within this paper's frame, be re-described and, in a sense, settled in Fuller's favour\,---\,not on Fuller's own grounds but on the paper's: the anchor of validity, reciprocity, is itself a normative and relational thing, a keeping of faith, and not a neutral fact about convergent behaviour; so the validity of law cannot after all be specified without reference to the relation of trust in which it is anchored. Fuller's inner morality is, in this paper's terms, the institutional-scale form of all three channels of \xref{p13:sec:mechanism} at once\,---\,a structure that is also a self-commitment and that invites a wager\,---\,with the anchor lodged in a doubly relied-upon promise.

\subsection{The anchor in procedure and discourse: Habermas and Rawls}
\label{p13:sub:jur-habermas}

A briefer station continues the movement Fuller began, by asking how legitimate norms are produced rather than where validity finally rests. Habermas locates the legitimacy of law in the discourse through which it is generated: norms are valid insofar as they could meet with the assent of all affected, in a communicative process governed by reason rather than force, and the legal-political system depends for its legitimacy upon a lifeworld of communicative practice that the system itself cannot produce and can only deplete~\parencite{habermas1996}. The structure of this dependence is worth marking, for it is isomorphic with a dependence this paper will meet again: just as the Habermasian system draws its legitimacy from a lifeworld it cannot generate from within itself, so the likelihood-channel of trust, in \xref{p13:sec:mechanism}, draws upon a real basis that the symbolic cannot manufacture. Rawls, with the device of a procedure\,---\,the original position\,---\,whose fairness confers legitimacy upon what issues from it, belongs to the same family~\parencite{rawls1971}: a family that anchors validity in the manner of its legitimate \emph{production}. This is structurally the same kind of question as the paper's own\,---\,how is something legitimately produced?\,---\,asked of legitimacy where the paper asks it of trust; and noting the kinship, the chapter returns to its main line.

\subsection{The first counter-coordinate: Holmes's bad man and legal realism}
\label{p13:sub:jur-holmes}

Against the movement toward acceptance and reciprocity, two positions hold the anchor down at force, and they must be given their due, for each marks a permanent coordinate the paper will need. The first is Holmes's. His advice to view the law through the eyes of the bad man\,---\,the man who cares nothing for the morality of the law and wishes only to know what the courts will in fact do to him, so that he may calculate\,---\,is the purest statement of the external, predictive stance~\parencite{holmes1897}. For the bad man, law just \emph{is} a prediction of official behaviour, and obligation just is the calculated probability of a consequence; he is the perfect evidentialist, the subject of pure mechanical cause, who recognises no internal point of view because he occupies none.

The paper's treatment of Holmes is twofold. On the one hand, the bad man is real and his stance is irreducible: he is the exact portrait of the evidentialist trust-grammar that \xref{p13:sec:translation} will have to take seriously, the subject for whom only the predictable consequence counts and for whom sincerity is no evidence. To pretend he does not exist, or that his stance is simply a moral failing, is to disarm the paper before its hardest case. On the other hand, his limit is precise and must be stated. To model \emph{every} subject as the bad man\,---\,to take the external predictive stance as the whole truth about legal subjects\,---\,is to erase the internal point of view a priori, and with it the possibility of meeting the other as a free subject who binds himself rather than as an object whose conduct is to be predicted and managed. The bad-man theory, generalised, is Kant's worry (\xref{p13:sub:diag-humekant}) realised in jurisprudence: it reduces the person to an object determined by consequences. It is true of those who have made themselves bad men, and it is a standing possibility for anyone; it is not the truth about the human as such, and a theory that takes it for that truth has decided in advance against everything \xref{p13:sec:freedom} will affirm. The sociological extension of realism\,---\,the \emph{living law} of Ehrlich, the law that actually governs conduct as against the law on the books\,---\,is noted here as a related and more capacious development, but the bad man is the figure that matters for the paper's argument.

\subsection{The second counter-coordinate: Kelsen's pure theory}
\label{p13:sub:jur-kelsen}

The second position that holds the anchor away from trust does so in the opposite direction from Holmes\,---\,not downward into predicted fact but upward into pure norm. Kelsen's pure theory derives the validity of each legal norm from a higher norm that authorises it, in an ascending chain that must terminate, on pain of regress, in a presupposed \emph{basic norm}, the \emph{Grundnorm}, which is not itself posited by any higher authority but assumed as the transcendental-logical condition of regarding the whole order as valid~\parencite{kelsen1967}. The Grundnorm is a presupposition, bracketing both morality and fact, a purely normative postulate at the apex of a purely normative structure.

Set beside Hart, the contrast is exact and instructive. Hart's rule of recognition is an accepted social fact, something officials actually do; Kelsen's Grundnorm is a presupposed symbolic posit, something one must assume. And the diagnosis the paper offers of Kelsen follows immediately: Kelsen lodges the anchor of validity in the purest possible point of the Symbolic, a posit suspended from nothing, held up by nothing but the decision to presuppose it. This is the most elegant of the positions and, by the paper's lights, the most fragile under the conditions the symptomatology described. For a foundation that consists in a pure symbolic presupposition is exactly what the crisis of the symbol dissolves: let the presupposition cease to be made\,---\,let people stop assuming the basic norm\,---\,and the whole towering normative structure, having no other support, hangs in the void. Kelsen is the limiting case of an anchor lodged in the Symbolic alone, and the crisis of trust is precisely the condition under which an anchor so lodged goes to zero (\xref{p13:sub:diag-trust}).

\subsection{The concrete instrument: from consideration to reliance}
\label{p13:sub:jur-estoppel}

The migration of the anchor is not only a movement in high theory; the law records it, with documentary precision, in the technical doctrines of contract, and these are the most directly relevant of all, for they bear on exactly the instrument from which the paper set out, the bare and unenforceable promise. The orthodox doctrine of \emph{consideration} holds that a bare promise\,---\,a promise for which nothing is given in exchange\,---\,is not enforceable; for a promise to bind, there must be a bargained-for exchange, a quid pro quo. The bare promise, the gratuitous undertaking, the memorandum that records an intention without exchanging anything\,---\,this is the precise legal situation of the betrothal document and of the memorandum of understanding alike: a promise the orthodox doctrine declines to enforce because nothing was given for it.

But the law developed an exception, and the exception is, for this paper, a quiet vindication from within the law itself. Under the doctrine of \emph{promissory estoppel}, a promise unsupported by consideration may nonetheless be enforced where the promisee has reasonably relied upon it, has acted to his detriment in that reliance, and would suffer injustice if the promise were not honoured. The ground of the binding force here is neither consideration nor any threat of sanction: it is \emph{reliance itself}. Because the promisee trusted the promise and made himself vulnerable in that trust, the promise acquires a binding force it did not have as a bare promise\,---\,a force generated by the trust placed in it, not by anything exchanged for it or any sword behind it. This is the mirror, in the law of contract, of the self-committing channel of \xref{p13:sec:mechanism}: the promisee's reliance is a placing of oneself in vulnerability, and the law responds to that vulnerability by generating an obligation to protect it. The associated measure, the \emph{reliance interest}\,---\,Fuller and Perdue's classic identification of the interest in being protected against the losses one suffers through having relied, as distinct from the expectation and restitution interests~\parencite{fullerperdue1936}\,---\,is the law's own name for the harm of a trust betrayed.

The significance is large and should be stated plainly. At its own technical edge, in the doctrine of promissory estoppel, the law has conceded that the true ground of a promise's binding force is reliance\,---\,is trust\,---\,and that enforcement is only the after-the-fact protection of a trust that was already, in itself and before any enforcement, the thing that bound. Promissory estoppel is the law's spontaneous, internal corroboration of this paper's central thesis: not a doctrine the paper imports to support its claim, but a place where the law, reasoning about its own most basic instrument, arrived at the claim by itself.

\subsection{The migration in one view}
\label{p13:sub:jur-summary}

The eight positions, read along the single axis of the anchor of validity, compose one directional movement. The law's own reflection upon itself has carried the anchor from force (Austin, Hobbes), through acceptance (Hart), to the reciprocal relation of reliance (Fuller, and the doctrine of promissory estoppel), with the predictive and the purely normative held at the two counter-coordinates (Holmes, Kelsen) that mark the limits the movement leaves behind. The paper does not impose this movement from outside; it traces an undercurrent already present in the law and connects it to the diagnosis of the symbol's crisis and, in \xref{p13:sec:mechanism}, to the mechanism of trust\,---\,giving the movement the full causal grounding the law itself, lacking a mechanistic vocabulary, could mark but not explain. The consequence is that this chapter is not mere background but an independent pillar of the argument: it establishes, from within the philosophy of law, that the anchor of validity is migrating from force toward trust, so that the metaphysics of \xref{p13:sec:freedom} and the mechanism of \xref{p13:sec:mechanism} are not free-standing philosophical or scientific assertions but the grounding explanation of a movement the law has already, on its own, undergone. Three lines\,---\,the law's internal undercurrent, the metaphysics of inferential cause, and the predictive-coding mechanism\,---\,converge on a single proposition: that it is trust, and not force, that is the true cause of obligation.

The whole may be set in one view. The table reads the positions along the axis of the chapter: where each lodges the anchor of validity, which sense of cause it thereby invokes, and what it contributes to, or marks the limit of, the paper's thesis.

\begin{center}
\small
\renewcommand{\arraystretch}{1.35}
\begin{longtable}{>{\RaggedRight\arraybackslash}p{2.9cm} >{\RaggedRight\arraybackslash}p{3.3cm} >{\RaggedRight\arraybackslash}p{2.3cm} >{\RaggedRight\arraybackslash}p{3.0cm}}
\toprule
\textcolor{rosedeep}{\textbf{\small Position}} & \textcolor{rosedeep}{\textbf{\small Anchor of validity}} & \textcolor{rosedeep}{\textbf{\small Sense of cause}} & \textcolor{rosedeep}{\textbf{\small Role in the thesis}} \\
\midrule
\endfirsthead
\toprule
\textcolor{rosedeep}{\textbf{\small Position}} & \textcolor{rosedeep}{\textbf{\small Anchor of validity}} & \textcolor{rosedeep}{\textbf{\small Sense of cause}} & \textcolor{rosedeep}{\textbf{\small Role in the thesis}} \\
\midrule
\endhead
\bottomrule
\endfoot
Austin, Hobbes & the threatened sanction; the sword & mechanical & states the position to be moved from; theory of a pathological limit \\
Hart & convergent acceptance; the internal point of view & inferential (normative) & the hinge: anchor pried loose from force toward acceptance \\
Fuller & reciprocity; mutual keeping of faith & inferential (relational) & names trust under its legal name; congruence as implicit promise \\
Habermas, Rawls & legitimate procedure; discourse & inferential (procedural) & isomorphic question: how is legitimacy produced \\
Holmes, realism & predicted official conduct & mechanical (predictive) & counter-coordinate; portrait of the evidentialist; its limit is Kant's worry \\
Kelsen & the presupposed basic norm & purely symbolic posit & counter-coordinate; anchor in the Symbolic alone; goes to zero in the crisis \\
Promissory estoppel & reliance itself & inferential (relational) & the law's own corroboration: reliance, not force, binds \\
\end{longtable}
\end{center}

\noindent
Read down the column of causes, the table shows the chapter's claim in a single glance: the entries that hold the anchor at force are entries of mechanical cause, and they sit at the beginning of the movement and at its abandoned limits; the entries that carry the anchor toward trust are entries of inferential cause, and they are where the law's reflection has been travelling. The migration of the anchor and the distinction between the two senses of cause are the same movement seen twice. What that inferential cause is, metaphysically\,---\,why it is the opening in which freedom itself becomes possible\,---\,is the question \xref{p13:sec:freedom} now takes up; how it works, mechanically, is the question of \xref{p13:sec:mechanism}.
